% GENERATED FILE. Edit canonical lessons, then run npm run generate:books.
% canonical-source-hash: fnv1a64:337a9897fc3e9233
% canonical-lessons: HI-C285-sabse, HI-C285-sezyada, HI-C285-dosaalse, HI-R285-comparing-and-time

\chapter{Taller, the Tallest, and How Long}
\label{ch:hi-taller-the-tallest-and-how-long}
\hlchaptermodality{285}

\begin{chapteropening}
\textbf{By the end of this chapter:} \emph{I can compare two things (se zyādā, se kam), say which is the most (sabse), and say how long something has gone on (do sāl se, din bhar).}

Complete HI-R285-comparing-and-time: compare, pick the most, and say how long.
\end{chapteropening}

\section[\texorpdfstring{\hi{सबसे} (sabse)}{सबसे (sabse)}]{\hi{सबसे} (sabse) — most, of all}

\label{lesson:HI-C285-sabse}



\begin{quote}
Before the new one: say \textquotedblleft{}all\textquotedblright{} (\textbf{\hi{सब}}) and \textquotedblleft{}the man is bigger than the child\textquotedblright{}. (\textbf{\hi{आदमी} \hi{बच्चे} \hi{से} \hi{बड़ा} \hi{है}}.)
\end{quote}

\subsection*{You'll want to know: \hi{सबसे}}
\textbf{\hi{सबसे}} — \emph{sabse} — \textquotedblleft{}most, of all\textquotedblright{}.

\begin{quote}
\textbf{\hi{मेरे} \hi{परिवार} \hi{में} \hi{मीरा} \hi{सबसे} \hi{लंबी} \hi{है।}}
\end{quote}

\begin{quote}
\emph{mere parivār meṁ Mīrā sabse lambī hai.} — \textquotedblleft{}In my family Mira is the tallest.\textquotedblright{}
\end{quote}

It is \textbf{\hi{सब}}, \textquotedblleft{}all\textquotedblright{}, with the \textbf{\hi{से}} of \textbf{\hi{बच्चे} \hi{से} \hi{बड़ा}}: \textquotedblleft{}taller than all\textquotedblright{}. The adjective after it agrees as always: \textbf{\hi{सबसे} \hi{लंबी}} for Mira, \textbf{\hi{सबसे} \hi{लंबा}} for Arun. The group you choose from takes \textbf{\hi{में}}.

\begin{itemize}
\raggedright
  \item \textbf{\hi{यह} \hi{सबसे} \hi{अच्छा} \hi{है।}} (\emph{yah sabse acchā hai.}) — This is the best.
  \item \textbf{\hi{सबसे} \hi{पहले}} (\emph{sabse pahle}) — first of all
\end{itemize}

\begin{cousinweb}
\textbf{\hi{सबसे}} (\emph{sabse}) is two words written as one: \textbf{\hi{सब}}, which CDIAL traces to Sanskrit \emph{sarva}, \textquotedblleft{}all, whole\textquotedblright{} (CDIAL 13276), and \textbf{\hi{से}}.
\end{cousinweb}

\subsection*{Your turn}
Say these aloud:

\begin{itemize}
\raggedright
  \item \emph{sabse}, then \emph{sabse acchā}
  \item \textquotedblleft{}In my family Mira is the tallest\textquotedblright{} — \emph{mere parivār meṁ Mīrā sabse lambī hai}
  \item \textquotedblleft{}First of all we drank tea\textquotedblright{} — \emph{sabse pahle hamne cāy pī}
\end{itemize}

\begin{tcolorbox}[breakable,colback=teal!4,colframe=teal!35!black,title={Before you move on}]
What does \textbf{\hi{सबसे}} mean? (\textquotedblleft{}Most, of all\textquotedblright{}.) \textquotedblleft{}The best\textquotedblright{}? (\textbf{\hi{सबसे} \hi{अच्छा}}.) \textquotedblleft{}First of all\textquotedblright{}? (\textbf{\hi{सबसे} \hi{पहले}}.)
\end{tcolorbox}
\section[\texorpdfstring{\hi{मीरा} \hi{अरुण} \hi{से} \hi{लंबी} \hi{है} (Mīrā Aruṇ se …)}{मीरा अरुण से लंबी है (Mīrā Aruṇ se …)}]{\hi{मीरा} \hi{अरुण} \hi{से} \hi{लंबी} \hi{है} (Mīrā Aruṇ se lambī hai) — comparing two, and choosing the most}

\label{lesson:HI-C285-sezyada}



\begin{quote}
Say \textquotedblleft{}the man is bigger than the child\textquotedblright{}, as you first learned it. (\textbf{\hi{आदमी} \hi{बच्चे} \hi{से} \hi{बड़ा} \hi{है}}.) Then \textquotedblleft{}the best\textquotedblright{}. (\textbf{\hi{सबसे} \hi{अच्छा}}.)
\end{quote}

\subsection*{You'll want to know: more and less}
\begin{quote}
\textbf{\hi{मीरा} \hi{अरुण} \hi{से} \hi{लंबी} \hi{है।}}
\end{quote}

\begin{quote}
\emph{Mīrā Aruṇ se lambī hai.} — \textquotedblleft{}Mira is taller than Arun.\textquotedblright{}
\end{quote}

You know the frame: what is measured against takes \textbf{\hi{से}}, and the adjective agrees with the one described, \textbf{\hi{लंबी}} for Mira. Two words sharpen it:

\begin{itemize}
\raggedright
  \item \textbf{\hi{से} \hi{ज़्यादा}} (\emph{se zyādā}), \textquotedblleft{}more than\textquotedblright{}: \textbf{\hi{यह} \hi{कमरा} \hi{उस} \hi{कमरे} \hi{से} \hi{ज़्यादा} \hi{बड़ा} \hi{है।}} (\emph{yah kamrā us kamre se zyādā baṛā hai.}) — This room is bigger than that one.
  \item \textbf{\hi{से} \hi{कम}} (\emph{se kam}), \textquotedblleft{}less than\textquotedblright{}: \textbf{\hi{यह} \hi{साड़ी} \hi{उससे} \hi{कम} \hi{महँगी} \hi{है।}} (\emph{yah sāṛī usse kam mahãgī hai.}) — This sari is less expensive than that one.
\end{itemize}

\textbf{\hi{उस} \hi{कमरे} \hi{से}} and \textbf{\hi{उससे}} both use \textbf{\hi{उस}}, the form \textbf{\hi{वह}} takes when a postposition follows, as in \textbf{\hi{उसके} \hi{बाद}}; \textbf{\hi{उससे}} is \textbf{\hi{उस}} with \textbf{\hi{से}}, \textquotedblleft{}than that one\textquotedblright{}.

They work with nouns too: \textbf{\hi{मेरे} \hi{पास} \hi{तुमसे} \hi{ज़्यादा} \hi{पैसे} \hi{हैं}} (\emph{mere pās tumse zyādā paise haiṁ}), \textquotedblleft{}I have more money than you\textquotedblright{}.

\subsection*{You'll want to know: the most, in a group}
\begin{quote}
\textbf{\hi{मुझे} \hi{फलों} \hi{में} \hi{आम} \hi{सबसे} \hi{ज़्यादा} \hi{पसंद} \hi{है।}}
\end{quote}

\begin{quote}
\emph{mujhe phaloṁ meṁ ām sabse zyādā pasand hai.} — \textquotedblleft{}Of all fruits I like mangoes the most.\textquotedblright{}
\end{quote}

\textbf{\hi{सबसे}} picks one from a group, and the group takes \textbf{\hi{में}}: \textbf{\hi{फलों} \hi{में}}, \textbf{\hi{इस} \hi{शहर} \hi{में}}. With a liking, \textquotedblleft{}the most\textquotedblright{} is \textbf{\hi{सबसे} \hi{ज़्यादा}}. A place can be the group too, with \textbf{\hi{का}}:

\begin{itemize}
\raggedright
  \item \textbf{\hi{यह} \hi{शहर} \hi{का} \hi{सबसे} \hi{बड़ा} \hi{बाज़ार} \hi{है।}} (\emph{yah śahar kā sabse baṛā bāzār hai.}) — This is the biggest market in the city.
\end{itemize}

\subsection*{Your turn}
Say these aloud:

\begin{itemize}
\raggedright
  \item \textquotedblleft{}This sari is less expensive than that one\textquotedblright{} — \emph{yah sāṛī usse kam mahãgī hai}
  \item \textquotedblleft{}I have more money than you\textquotedblright{} — \emph{mere pās tumse zyādā paise haiṁ}
  \item \textquotedblleft{}Of all fruits I like mangoes the most\textquotedblright{} — \emph{mujhe phaloṁ meṁ ām sabse zyādā pasand hai}
\end{itemize}

\begin{tcolorbox}[breakable,colback=teal!4,colframe=teal!35!black,title={Before you move on}]
\textquotedblleft{}Taller than Arun\textquotedblright{}? (\textbf{\hi{अरुण} \hi{से} \hi{लंबी}}, or \textbf{\hi{लंबा}}.) \textquotedblleft{}More than\textquotedblright{}? (\textbf{\hi{से} \hi{ज़्यादा}}.) \textquotedblleft{}Less than\textquotedblright{}? (\textbf{\hi{से} \hi{कम}}.) \textquotedblleft{}The biggest market in the city\textquotedblright{}? (\textbf{\hi{शहर} \hi{का} \hi{सबसे} \hi{बड़ा} \hi{बाज़ार}}.)
\end{tcolorbox}
\section[\texorpdfstring{\hi{दो} \hi{साल} \hi{से} (do sāl se)}{दो साल से (do sāl se)}]{\hi{दो} \hi{साल} \hi{से} (do sāl se) — for how long}

\label{lesson:HI-C285-dosaalse}



\begin{quote}
Say \textquotedblleft{}from morning to evening\textquotedblright{}. (\textbf{\hi{सुबह} \hi{से} \hi{शाम} \hi{तक}}.) Then \textquotedblleft{}I am doing\textquotedblright{}. (\textbf{\hi{मैं} \hi{कर} \hi{रहा} \hi{हूँ}}.)
\end{quote}

\subsection*{You'll want to know: since, and for}
\begin{quote}
\textbf{\hi{मैं} \hi{दो} \hi{साल} \hi{से} \hi{हिंदी} \hi{सीख} \hi{रहा} \hi{हूँ।}}
\end{quote}

\begin{quote}
\emph{maiṁ do sāl se hindī sīkh rahā hūṁ.} — \textquotedblleft{}I have been learning Hindi for two years.\textquotedblright{}
\end{quote}

English says \textquotedblleft{}have been learning\textquotedblright{}. Hindi says \textquotedblleft{}am learning\textquotedblright{}, in the present, because it is still going on, and puts \textbf{\hi{से}}, \textquotedblleft{}from\textquotedblright{}, on the stretch so far or on the time it began:

\begin{itemize}
\raggedright
  \item \textbf{\hi{बारिश} \hi{सुबह} \hi{से} \hi{हो} \hi{रही} \hi{है।}} (\emph{bāriś subah se ho rahī hai.}) — It has been raining since the morning.
  \item \textbf{\hi{वह} \hi{पाँच} \hi{साल} \hi{से} \hi{यहाँ} \hi{रहती} \hi{है।}} (\emph{vah pā̃c sāl se yahā̃ rahtī hai.}) — She has lived here for five years.
  \item \textbf{\hi{कब} \hi{से}?} (\emph{kab se?}) — Since when? \textbf{\hi{कितने} \hi{दिनों} \hi{से}?} (\emph{kitne dinoṁ se?}) — For how many days?
\end{itemize}

\textbf{\hi{से} … \hi{तक}}, which you know, gives a stretch with both ends: \textbf{\hi{सुबह} \hi{से} \hi{शाम} \hi{तक}}. \textbf{\hi{से}} on its own counts up to now.

\subsection*{You'll want to know: all day long}
\begin{quote}
\textbf{\hi{कल} \hi{मैं} \hi{दिन} \hi{भर} \hi{काम} \hi{कर} \hi{रहा} \hi{था।}}
\end{quote}

\begin{quote}
\emph{kal maiṁ din bhar kām kar rahā thā.} — \textquotedblleft{}Yesterday I was working all day.\textquotedblright{}
\end{quote}

\textbf{\hi{भर}} (\emph{bhar}) after a stretch of time is \textquotedblleft{}all through it\textquotedblright{}: \textbf{\hi{दिन} \hi{भर}} (\emph{din bhar}), all day; \textbf{\hi{रात} \hi{भर}} (\emph{rāt bhar}), all night; \textbf{\hi{साल} \hi{भर}} (\emph{sāl bhar}), all year. \textbf{\hi{पूरा}} says the same as an adjective: \textbf{\hi{पूरे} \hi{दिन}} (\emph{pūre din}), the whole day.

\begin{itemize}
\raggedright
  \item \textbf{\hi{रात} \hi{भर} \hi{बारिश} \hi{हो} \hi{रही} \hi{थी।}} (\emph{rāt bhar bāriś ho rahī thī.}) — It was raining all night.
\end{itemize}

\subsection*{Your turn}
Say these aloud:

\begin{itemize}
\raggedright
  \item \textquotedblleft{}Since when?\textquotedblright{} — \emph{kab se?}
  \item \textquotedblleft{}It has been raining since the morning\textquotedblright{} — \emph{bāriś subah se ho rahī hai}
  \item \textquotedblleft{}She has lived here for five years\textquotedblright{} — \emph{vah pā̃c sāl se yahā̃ rahtī hai}
\end{itemize}

\begin{tcolorbox}[breakable,colback=teal!4,colframe=teal!35!black,title={Before you move on}]
\textquotedblleft{}For two years\textquotedblright{}? (\textbf{\hi{दो} \hi{साल} \hi{से}}.) Why the present in \textbf{\hi{मैं} \hi{दो} \hi{साल} \hi{से} \hi{हिंदी} \hi{सीख} \hi{रहा} \hi{हूँ}}? (It is still going on.) \textquotedblleft{}All night\textquotedblright{}? (\textbf{\hi{रात} \hi{भर}}.) \textquotedblleft{}From morning to evening\textquotedblright{}? (\textbf{\hi{सुबह} \hi{से} \hi{शाम} \hi{तक}}.)
\end{tcolorbox}
\section[(review)]{Review — comparing, and how long}

\label{lesson:HI-R285-comparing-and-time}



\begin{quote}
Say \textquotedblleft{}Mira is taller than Arun\textquotedblright{}. (\emph{Mīrā Aruṇ se lambī hai}.) Then \textquotedblleft{}the best\textquotedblright{}. (\emph{sabse acchā}.)
\end{quote}

\subsection*{Your turn}
Say these aloud:

\begin{itemize}
\raggedright
  \item This shop is more expensive than that one.
  \item In my opinion this is the best book in the library.
  \item I have wanted to become a teacher for two years.
  \item Yesterday it rained all day.
  \item Since when have you lived here?
\end{itemize}

\begin{tcolorbox}[breakable,colback=teal!4,colframe=teal!35!black,title={Before you move on}]
\textquotedblleft{}More expensive than\textquotedblright{}? (\textbf{\hi{से} \hi{ज़्यादा} \hi{महँगा}}.) \textquotedblleft{}The best in the library\textquotedblright{}? (\textbf{\hi{पुस्तकालय} \hi{की} \hi{सबसे} \hi{अच्छी} \hi{किताब}}: \textbf{\hi{किताब}} is feminine.) \textquotedblleft{}For two years\textquotedblright{}? (\textbf{\hi{दो} \hi{साल} \hi{से}}, with the present.) \textquotedblleft{}All day\textquotedblright{}? (\textbf{\hi{दिन} \hi{भर}}.)
\end{tcolorbox}
